% GENERATED FILE. Edit canonical lessons, then run npm run generate:books.
% canonical-source-hash: fnv1a64:bc81aa718c33e504
% canonical-lessons: ES-C357-navidad, ES-C357-pascua, ES-C357-vispera, ES-C357-regalo, ES-C357-tarjeta

\chapter{The Days A Calendar Circles}
\label{ch:la-navidad-y-la-pascua}
\hlchaptermodality{357}

\begin{chapteropening}
\textbf{By the end of this chapter:} \emph{I can say la Navidad, la Pascua, la vispera, el regalo and la tarjeta, and name the two days a Spanish calendar circles first.}

Complete ES-C357-tarjeta: name a feast, name the day before it, and name the two things that arrive with it.
\end{chapteropening}

\section[la Navidad]{la Navidad — the birth, and the nothing that shares its root}

\label{lesson:ES-C357-navidad}



\begin{quote}
Say \emph{the occasion}, then \emph{suddenly}. (\emph{La ocasión}; \emph{de repente}.)
\end{quote}

\subsection*{What to know first}
\begin{itemize}
\raggedright
  \item nada — the same Latin verb, and the opposite meaning.
  \item nadie — and the same verb again.
\end{itemize}

\begin{sounds}
\begin{itemize}
\raggedright
  \item \emph{na-vi-DAD}, three beats with the weight on the last. The \emph{v} sounds like a soft \emph{b}. $\to$ pronunciation reference
\end{itemize}
\end{sounds}

\begin{cousinweb}
\textbf{la Navidad} — \textquotedblleft{}Christmas.\textquotedblright{}

Latin \textbf{nāscī} is to be born. \textbf{Nātīvitātem} is a birth. And \textbf{la Navidad} is \emph{the} birth, with no need to say whose — the way \emph{la Guerra} in a Spanish sentence about 1936 needs no adjective.

English kept the Latin whole in \textbf{nativity}, and the root in \textbf{native} (born here), \textbf{nation} (people of one birth), \textbf{nature} (what a thing is born as), and \textbf{innate}.

Now the part worth stopping for. Two of the commonest words you own are the same verb, and they mean the opposite:

\begin{itemize}
\raggedright
  \item \textbf{nada} is Latin \emph{(rēs) nāta}, \textquotedblleft{}a thing born.\textquotedblright{} Latin said \emph{nōn... rem nātam}, \textquotedblleft{}not a born thing,\textquotedblright{} and Spanish threw away the \emph{non} and kept the \emph{nata}, so the word for \emph{nothing} is literally \emph{a born thing};
  \item \textbf{nadie} is \emph{(hominēs) nātī}, \textquotedblleft{}born people,\textquotedblright{} and did the same disappearing act with its negative.
\end{itemize}

So \emph{nada}, \emph{nadie} and \emph{la Navidad} are one Latin verb. Three births, and two of them ended up meaning nobody and nothing.
\end{cousinweb}

\begin{grammarlens}[title={What you've built}]
The first of the two days a Spanish calendar circles, and a root you have been using without knowing it.
\end{grammarlens}

\subsection*{Your turn}
Say these aloud:

\begin{itemize}
\raggedright
  \item \emph{la Navidad}
  \item \emph{en Navidad}
  \item \emph{feliz Navidad}
\end{itemize}

\begin{tcolorbox}[breakable,colback=teal!4,colframe=teal!35!black,title={Before you move on}]
What does \emph{nāscī} mean? (To be born.) And what does \emph{nada} mean word for word? (A thing born.)
\end{tcolorbox}
\section[la Pascua]{la Pascua — a Hebrew word that crossed three languages}

\label{lesson:ES-C357-pascua}



\begin{quote}
Say \emph{suddenly}, then \emph{Christmas}. (\emph{De repente}; \emph{la Navidad}.)
\end{quote}

\subsection*{What to know first}
\begin{itemize}
\raggedright
  \item la fiesta — what it is.
  \item la Navidad — the other one.
\end{itemize}

\begin{sounds}
\begin{itemize}
\raggedright
  \item \emph{PAS-kua}, two beats with the weight on the first. The \emph{cu} is a hard \emph{k} with a \emph{w} glide. $\to$ pronunciation reference
\end{itemize}
\end{sounds}

\begin{cousinweb}
\textbf{la Pascua} — \textquotedblleft{}Easter.\textquotedblright{}

This word has crossed more borders than most. Hebrew \textbf{pesach} is the Passover. Greek wrote it \textbf{páskha}, Latin \textbf{pascha}, and from Latin it went into every Romance language: Spanish \textbf{Pascua}, Italian \emph{Pasqua}, French \emph{Pâques}, Portuguese \emph{Páscoa}.

English is the odd one out, and this is the interesting part. For the feast itself English kept a Germanic name — \textbf{Easter}, which Bede says came from a spring goddess — but it kept the Latin word for the adjective. So English has the \textbf{paschal} lamb and the \textbf{paschal} moon, and it has Easter eggs, and the two words never meet.

The \emph{u} in \textbf{Pascua} is a passenger. Latin also had \textbf{pascua}, \textquotedblleft{}pastures,\textquotedblright{} and the two words sat close enough for centuries that the pasture lent its shape to the feast — helped along by the paschal lamb, which grazes.

Spanish also uses \emph{Pascua} in the plural for the Christmas season, so \emph{felices Pascuas} can mean the wrong feast to an English ear.
\end{cousinweb}

\begin{grammarlens}[title={What you've built}]
Both of the days a Spanish year turns on, and one Hebrew word underneath one of them.
\end{grammarlens}

\subsection*{Your turn}
Say these aloud:

\begin{itemize}
\raggedright
  \item \emph{la Pascua}
  \item \emph{en Pascua}
  \item \emph{la Navidad y la Pascua}
\end{itemize}

\begin{tcolorbox}[breakable,colback=teal!4,colframe=teal!35!black,title={Before you move on}]
Which Hebrew festival is \emph{Pascua} named after? (The Passover, \emph{pesach}.) And which English adjective kept the Latin word? (Paschal.)
\end{tcolorbox}
\section[la víspera]{la víspera — the evening that leans on the next day}

\label{lesson:ES-C357-vispera}



\begin{quote}
Say \emph{Christmas}, then \emph{Easter}. (\emph{La Navidad}; \emph{la Pascua}.)
\end{quote}

\subsection*{What to know first}
\begin{itemize}
\raggedright
  \item el oeste — the direction the evening is in.
  \item antes — what it means, reduced to one word.
\end{itemize}

\begin{sounds}
\begin{itemize}
\raggedright
  \item \emph{VIS-pe-ra}, three beats with the weight on the FIRST, which is unusual enough that it carries a written accent. $\to$ pronunciation reference
\end{itemize}
\end{sounds}

\begin{cousinweb}
\textbf{la víspera} — \textquotedblleft{}the eve,\textquotedblright{} the day before.

Latin \textbf{vespera} is the evening. Spanish narrowed it to a particular evening: not any evening, but the one that leans on a bigger day. \emph{La víspera de Navidad} is Christmas Eve. \emph{La víspera del viaje} is the night before you leave.

English has the word in one place, and it is a quiet one. \textbf{Vespers} are the evening prayers, sung as the light goes.

Greek had the same ancient word with the \emph{v} worn off the front: \textbf{hésperos}, the evening star, which is Venus seen after sunset. The \textbf{Hesperides}, whose garden held the golden apples, are the daughters of evening, and their garden was at the far western edge of the world — because that is where the evening is.

Which is the last piece. English \textbf{west} is a cousin of \emph{vespera} and \emph{hesperos}, all three from an ancient root for \emph{down} or \emph{late}. So \emph{el oeste} and \emph{la víspera} point at the same thing from two sides: the place the sun goes down, and the hours after it does.
\end{cousinweb}

\begin{grammarlens}[title={What you've built}]
A way to name the day before a day, without counting backwards.
\end{grammarlens}

\subsection*{Your turn}
Say these aloud:

\begin{itemize}
\raggedright
  \item \emph{la víspera}
  \item \emph{la víspera de Navidad}
  \item \emph{la víspera de la fiesta}
\end{itemize}

\begin{tcolorbox}[breakable,colback=teal!4,colframe=teal!35!black,title={Before you move on}]
What does \emph{vespera} mean? (Evening.) And which English direction is a cousin of it? (West.)
\end{tcolorbox}
\section[el regalo]{el regalo — merry-making that turned into an object}

\label{lesson:ES-C357-regalo}



\begin{quote}
Say \emph{Easter}, then \emph{the eve}. (\emph{La Pascua}; \emph{la víspera}.)
\end{quote}

\subsection*{What to know first}
\begin{itemize}
\raggedright
  \item la fiesta — when it arrives.
  \item el placer — what it is for.
\end{itemize}

\begin{sounds}
\begin{itemize}
\raggedright
  \item \emph{rre-GA-lo}, three beats with the weight on the second. The opening \emph{r} is rolled. The \emph{g} before \emph{a} is hard. $\to$ pronunciation reference
\end{itemize}
\end{sounds}

\begin{cousinweb}
\textbf{el regalo} — \textquotedblleft{}the present,\textquotedblright{} \textquotedblleft{}the gift.\textquotedblright{}

Spanish borrowed the verb \textbf{regalar} from French \emph{régaler}, to entertain somebody handsomely — to treat them. Under it is an Old French noun, \textbf{gale}, meaning merry-making or pleasure.

English took the same French verb and kept the entertaining: to \textbf{regale} a guest is to give them a fine story or a fine dinner. And \emph{gale} on its own gave English \textbf{gala}, a day of merry-making.

Then Spanish did something the other languages did not. It let the word slide from the entertaining to the \emph{object handed over}. In French you \emph{régale} someone; in Spanish you \emph{regalas} them a thing, and the thing itself is \textbf{un regalo}.

One warning, because the shape invites it. \textbf{Regalo} looks like \emph{real} and \emph{regio}, which come from Latin \textbf{rēx}, a king, and give English \textbf{regal} and \textbf{royal}. There is no connection. A \emph{regalo} is about a good time, not about a crown.
\end{cousinweb}

\begin{grammarlens}[title={What you've built}]
The thing that arrives on a feast day, and a false lead refused.
\end{grammarlens}

\subsection*{Your turn}
Say these aloud:

\begin{itemize}
\raggedright
  \item \emph{el regalo}
  \item \emph{un regalo}
  \item \emph{un regalo de Navidad}
\end{itemize}

\begin{tcolorbox}[breakable,colback=teal!4,colframe=teal!35!black,title={Before you move on}]
What did Old French \emph{gale} mean? (Merry-making.) And is \emph{regalo} related to \emph{real}, a king? (No — it only looks like it.)
\end{tcolorbox}
\section[la tarjeta]{la tarjeta — a small shield you write on}

\label{lesson:ES-C357-tarjeta}



\begin{quote}
Say \emph{the eve}, then \emph{the present}. (\emph{La víspera}; \emph{el regalo}.)
\end{quote}

\subsection*{What to know first}
\begin{itemize}
\raggedright
  \item la carta — its longer cousin.
  \item el regalo — what it travels with.
\end{itemize}

\begin{sounds}
\begin{itemize}
\raggedright
  \item \emph{tar-JE-ta}, three beats with the weight on the second. The \emph{j} is the rough sound at the back of the mouth. $\to$ pronunciation reference
\end{itemize}
\end{sounds}

\begin{cousinweb}
\textbf{la tarjeta} — \textquotedblleft{}the card.\textquotedblright{}

A \emph{tarjeta} is a little \textbf{tarja}, and a \emph{tarja} was a shield. Spanish took it from French \textbf{targe}, and French from a Frankish \textbf{targa} — a small round shield carried on the arm.

English took the same shield and made it small in its own way. \textbf{Target} is \emph{targe} plus a diminutive ending, and for several hundred years it meant exactly that: a light shield. Only later did somebody set one up in a field to shoot at, and the name went with the thing being aimed at rather than the thing doing the protecting.

What carried the Spanish word from armour to stationery is the shape. A shield is a small stiff rectangle, or something near it, held up for other people to read: a coat of arms is a shield that tells you who is behind it.

A \textbf{tarjeta} does the same job with ink. \emph{Una tarjeta de visita} announces who has arrived, and \emph{una tarjeta de Navidad} says who remembered you.
\end{cousinweb}

\begin{grammarlens}[title={What you've built}]
Five things a calendar day brings with it: \emph{la Navidad}, \emph{la Pascua}, \emph{la víspera}, \emph{el regalo}, \emph{la tarjeta}.
\end{grammarlens}

\subsection*{Your turn}
Say these aloud:

\begin{itemize}
\raggedright
  \item \emph{la tarjeta}
  \item \emph{una tarjeta y un regalo}
  \item \emph{la Navidad, la Pascua, la víspera}
\end{itemize}

\begin{tcolorbox}[breakable,colback=teal!4,colframe=teal!35!black,title={Before you move on}]
What was a \emph{targe}? (A small shield.) And what did \textbf{target} mean before anybody shot at one? (A shield.)
\end{tcolorbox}
